% A small document that exercises the parts of a LaTeX preview worth looking at:
% sectioning, display and inline math, a table, a figure-less float, and a
% cross-reference, which is the thing that needs a second compilation pass.
\documentclass[11pt]{article}

\usepackage{amsmath}
\usepackage{booktabs}

\title{MarkCopy LaTeX Preview}
\author{Sample Document}
\date{}

\begin{document}
\maketitle

\section{Inline and display math}
\label{sec:math}

Euler's identity, $e^{i\pi} + 1 = 0$, is compact enough to sit inline. The
Gaussian integral wants its own line:

\begin{equation}
  \label{eq:gauss}
  \int_{-\infty}^{\infty} e^{-x^2} \, dx = \sqrt{\pi}.
\end{equation}

\section{A cross-reference}

Equation~\eqref{eq:gauss} appears in Section~\ref{sec:math}. A reference like
this one resolves only on the second pass, which is why the preview prefers an
engine that reruns the document on its own.

\section{A table}

\begin{table}[h]
  \centering
  \begin{tabular}{lrr}
    \toprule
    Engine    & Passes & Fetches packages \\
    \midrule
    pdflatex  & manual & no  \\
    latexmk   & auto   & no  \\
    tectonic  & auto   & yes \\
    \bottomrule
  \end{tabular}
  \caption{What the preview is choosing between.}
\end{table}

\end{document}
